% ============================================================================
% STUNDENLEISTUNG (Papier): parecer y conocer
% ============================================================================
% 25 BE | 14-NP-Skala | Einzelarbeit
% ============================================================================

\documentclass[11pt,a4paper]{article}
\usepackage[utf8]{inputenc}
\usepackage[T1]{fontenc}
\usepackage[ngerman]{babel}
\usepackage{geometry}
\usepackage{fancyhdr}
\usepackage{xcolor}
\usepackage{enumitem}
\usepackage{tabularx}
\usepackage{array}

\geometry{left=2cm, right=2cm, top=2cm, bottom=2cm}
\definecolor{fach}{HTML}{c41e3a}
\setlength{\parindent}{0pt}

\pagestyle{fancy}
\fancyhf{}
\fancyhead[L]{\textbf{Stundenleistung: parecer y conocer}}
\fancyhead[R]{Spanisch Klasse 9}
\fancyfoot[C]{\small\thepage}
\renewcommand{\headrulewidth}{0.4pt}

\newcommand{\luecke}[1]{\underline{\hspace{#1}}}
\newcommand{\punktebox}[1]{\hfill\fbox{\small \_/#1}}

\begin{document}

% Kopfbereich
\begin{center}
{\Large\textcolor{fach}{\textbf{Stundenleistung: parecer y conocer}}}
\end{center}

\vspace{0.3cm}

\begin{tabular}{p{8cm}p{6cm}}
Name: \luecke{5cm} & Datum: \luecke{3cm} \\
\end{tabular}

\vspace{0.5cm}

% ============================================================================
% AUFGABE 1
% ============================================================================
\textbf{Aufgabe 1: Ergänze die richtige Form von \textit{parecer}.} \punktebox{6}

\begin{enumerate}[label=\alph)]
    \item ¿Qué te \luecke{2.5cm} los deberes?
    \item Me \luecke{2.5cm} muy difíciles.
    \item ¿Qué os \luecke{2.5cm} la película?
    \item Nos \luecke{2.5cm} muy divertida.
    \item A Clara le \luecke{2.5cm} que Gonzalo es raro.
    \item ¿Qué te \luecke{2.5cm} quedar a las ocho?
\end{enumerate}

\vspace{0.5cm}

% ============================================================================
% AUFGABE 2
% ============================================================================
\textbf{Aufgabe 2: Ergänze die richtige Form von \textit{conocer}.} \punktebox{7}

\begin{enumerate}[label=\alph)]
    \item Yo no \luecke{2.5cm} a tu hermano.
    \item ¿Tú \luecke{2.5cm} la calle Trinidad?
    \item Nosotros \luecke{2.5cm} a toda tu familia.
    \item Mis hermanos \luecke{2.5cm} a Alfonso.
    \item Clara todavía no \luecke{2.5cm} muy bien a Nicolás.
    \item ¿Vosotros \luecke{2.5cm} a los amigos de Gonzalo?
    \item Ado dice: ``Yo \luecke{2.5cm} todas las calles del centro.''
\end{enumerate}

\vspace{0.5cm}

% ============================================================================
% AUFGABE 3
% ============================================================================
\textbf{Aufgabe 3: Ergänze \textit{a} wo nötig. Wenn kein \textit{a} nötig ist, schreibe ``--''.} \punktebox{6}

\begin{enumerate}[label=\alph)]
    \item ¿Conoces \luecke{1cm} mi madre?
    \item No conozco \luecke{1cm} esta ciudad.
    \item Queremos conocer \luecke{1cm} tus amigos.
    \item ¿Conocéis \luecke{1cm} el nuevo piso de Clara?
    \item Todavía no conozco \luecke{1cm} Nicolás.
    \item Mi hermano conoce \luecke{1cm} todas las canciones.
\end{enumerate}

\vspace{0.5cm}

% ============================================================================
% AUFGABE 4
% ============================================================================
\textbf{Aufgabe 4: Kreuze die richtige Antwort an.} \punktebox{6}

\textbf{a)} ``Me parece simpático'' bedeutet: \hfill (2 P.)

$\Box$ Ich mag ihn. \hspace{1cm} $\Box$ Ich finde ihn sympathisch. \hspace{1cm} $\Box$ Er scheint krank.

\vspace{0.4cm}

\textbf{b)} Die yo-Form von ``conocer'' ist: \hfill (2 P.)

$\Box$ conoco \hspace{1cm} $\Box$ conozco \hspace{1cm} $\Box$ conosco

\vspace{0.4cm}

\textbf{c)} Wann steht ``a'' nach conocer? \hfill (2 P.)

$\Box$ Bei Sachen \hspace{1cm} $\Box$ Bei Personen \hspace{1cm} $\Box$ Immer

\vspace{1cm}

% ============================================================================
% NOTENSPIEGEL
% ============================================================================
\hrule
\vspace{0.3cm}

\begin{center}
\begin{tabular}{|c|c|c|c|c|c|c|c|c|c|c|c|c|c|c|c|}
\hline
\textbf{NP} & 14 & 13 & 12 & 11 & 10 & 9 & 8 & 7 & 6 & 5 & 4 & 3 & 2 & 1 & 0 \\
\hline
\textbf{Punkte} & 25 & 24 & 23 & 22 & 20 & 18 & 17 & 15 & 13 & 12 & 10 & 8 & 7 & 5 & $<$5 \\
\hline
\end{tabular}
\end{center}

\vspace{0.5cm}

\begin{flushright}
\textbf{Punkte:} \luecke{1.5cm} / 25 \hspace{1cm} \textbf{Note:} \luecke{1.5cm}
\end{flushright}

\end{document}
